\documentclass[12pt]{evaluasep}
\grado{3$^\circ$ de Secundaria}
\cicloescolar{2025-2026}
\materia{Olimpiada de Conocimientos 2026}
\unidad{1}
\title{Propuesta para Examen}
\author{Profesores de Secundaria}
\begin{document}
% \begin{multicols}{2}
	\tableofcontents
% \end{multicols}%
\newpage
\begin{questions}
		\section{Lenjuages}
                \subsection{Lengua Materna}

                \question ¿Cuál es el primer paso metodológico indispensable que se debe llevar a cabo para elaborar un texto informativo estructurado?
    
                \begin{choices}
        \choice Escribir el borrador definitivo sin necesidad de realizar lecturas previas o esquemas.
        \choice Diseñar una portada llamativa con recursos tipográficos y logotipos institucionales.
        \CorrectChoice Seleccionar un tema de interés y recopilar información en fuentes confiables y documentadas.
        \choice Publicar el escrito en medios digitales y redes sociales para recibir retroalimentación del público.
    
    \end{choices}
    

    \question ¿Qué característica distingue de manera contundente a una campaña publicitaria frente a una campaña de carácter social?
    
    \begin{choices}
        \choice Emplea mensajes objetivos validados científicamente y evita el uso de recursos visuales exagerados.
        \choice Está promovida de forma exclusiva por organismos internacionales y fundaciones comunitarias sin fines de lucro.
        \CorrectChoice Su meta primordial es persuadir al consumidor de los beneficios de un producto o servicio para incentivar su compra o contratación comercial.
        \choice Busca transformar de manera altruista los hábitos perjudiciales de una población vulnerable sin realizar transacciones monetarias.
    
    \end{choices}
    

    \question Si se diseña una campaña social con un cartel enfocado en mitigar los padecimientos crónicos derivados del consumo de bebidas azucaradas, ¿cuál es su formato y su eje temático?
    
    \begin{choices}
        \choice El formato es un folleto informativo de tres caras y el tema principal es el cuidado integral del agua potable.
        \CorrectChoice El formato utilizado es un cartel de impacto visual directo y el tema de estudio es el cuidado de la salud.
        \choice El formato elegido es una infografía secuencial y el tema central es la protección de menores de edad.
        \choice El formato es una revista comunitaria y el tema trata sobre la producción sustentable de alimentos orgánicos.
    
    \end{choices}
    

    \question En el contexto escolar, ¿qué ventaja representan los medios de comunicación impresos tales como carteles y gacetas?
    
    \begin{choices}
        \choice Son la única alternativa tecnológica que garantiza el cumplimiento total de los objetivos de la institución educativa.
        \CorrectChoice Son una excelente vía para asegurar que la comunidad del plantel lea, analice y conozca los pormenores de una campaña de difusión interna.
        \choice Constituyen el método más rápido, moderno e infalible para resolver las problemáticas administrativas dentro de la escuela.
        \choice Reemplazan de forma definitiva las sesiones académicas presenciales mediante la entrega autónoma de información.
    
    \end{choices}
    
\newpage
    \question Ante un conflicto en el entorno escolar, ¿cuál se considera la vía más eficiente y formativa para alcanzar una resolución pacífica?
    
    \begin{choices}
        \choice Decretar sanciones punitivas muy severas y reportes disciplinarios directos para todas las partes involucradas.
        \choice Diseñar materiales impresos de amonestación pública y distribuirlos por las aulas del plantel escolar.
        \CorrectChoice Exponer abiertamente la problemática detectada y construir soluciones viables mediante el trabajo colaborativo.
        \choice Evitar el diálogo directo y delegar toda la responsabilidad en manos de las autoridades externas al plantel.
    
    \end{choices}
    

    \question ¿Qué acción de control y cierre técnico se debe ejecutar obligatoriamente al finalizar cualquier tipo de campaña de difusión masiva?
    
    \begin{choices}
        \choice Organizar una exposición artística pública que recopile exclusivamente los mejores diseños visuales de los alumnos.
        \choice Colocar un buzón de sugerencias permanente en la entrada de las instalaciones para recibir quejas generales.
        \CorrectChoice Implementar una evaluación metodológica rigurosa con el fin de medir la repercusión y el impacto real del mensaje difundido.
        \choice Destruir de inmediato todos los materiales físicos generados para evitar la acumulación de desechos en las aulas.
    
    \end{choices}
    

    \question ¿Cuál de las siguientes alternativas representa un desarrollo fundamental del conocimiento humano que NO deriva del método científico moderno?
    
    \begin{choices}
        \choice El diseño, síntesis y dosificación segura de medicamentos avanzados.
        \choice El procesamiento químico industrial para la conservación a largo plazo de comida enlatada.
        \CorrectChoice El origen, diseño evolutivo y diversificación de los sistemas de escritura.
        \choice El desarrollo tecnológico de microprocesadores para la fabricación de dispositivos informáticos.
    
    \end{choices}
    

    \question ¿Cuál de los siguientes enunciados define con exactitud el objeto de estudio de los textos científicos?
    
    \begin{choices}
        \choice El análisis ontológico del pensamiento racional y las corrientes de la filosofía clásica.
        \CorrectChoice La explicación fundamentada del funcionamiento intrínseco de la naturaleza y sus leyes físicas.
        \choice La recopilación e interpretación mitológica de las leyendas y los orígenes culturales de las civilizaciones antiguas.
        \choice La creación y el perfeccionamiento técnico de obras narrativas basadas puramente en la ficción literaria.
    
    \end{choices}
    

                \question ¿Cuál es la función de las convocatorias?
    
                \begin{choices}
        \choice Expresar sentimientos íntimos del autor a través de versos y estrofas.
        \CorrectChoice Invitar a participar en un evento o concurso.
        \choice Informar de manera objetiva sobre un acontecimiento reciente.
        \choice Entablar una conversación informal entre dos amigos lejanos.
    
    \end{choices}
                
\newpage
    \question ¿Cuáles son dos características de los textos literarios?
    
    \begin{choices}
        \choice Cuentan con un título, introducción, desarrollo, conclusión y fuentes de consulta.
        \CorrectChoice Se emplean en cartas y funcionan como una técnica de discusión donde se exponen ideas y puntos de vista fundamentados.
        \choice Son textos breves que narran una historia de principio a fin de forma muy rápida.
        \choice Su única finalidad es convocar a personas a un evento público.
    
    \end{choices}
    

    \question ¿Cuáles son las partes que conforman un artículo informativo?
    
    \begin{choices}
        \choice Inicio, nudo, clímax y desenlace.
        \choice Remitente, destinatario, saludo y despedida.
        \CorrectChoice Título, introducción, desarrollo, conclusión y fuentes de consulta.
        \choice Verso, estrofa, rima y métrica.
    
    \end{choices}
    

    \question ¿A qué se le llama micorrelatos?
    
    \begin{choices}
        \choice A las obras teatrales divididas en tres actos o jornadas.
        \choice A los artículos de opinión publicados en los periódicos.
        \CorrectChoice A los textos narrativos que cuentan una historia de principio a fin.
        \choice A los textos formales que se envían a un destinatario específico.
    
    \end{choices}
    

    \question ¿Qué son las cartas o epístolas?
    
    \begin{choices}
        \choice Textos que sirven exclusivamente para invitar a concursos.
        \CorrectChoice Son textos formales o personales en los que un remitente entabla comunicación con un destinatario específico.
        \choice Textos narrativos breves que contienen un inicio y un desenlace rápido.
        \choice Artículos cuyo objetivo es mostrar fuentes de consulta sobre un tema.
    
    \end{choices}
    

    \subsection{Lengua extranjera}
    \question What presents did you \underline{\hspace{2cm}} for your last birthday?
    
    \begin{choices}
        \choice will receive
        \CorrectChoice receive
        \choice receiving
        \choice received
    
    \end{choices}
    

    \question We use ``someone'', ``anyone'' and ``no one'' to talk about:
    
    \begin{choices}
        \CorrectChoice people
        \choice locations
        \choice things
        \choice numbers
    
    \end{choices}
    
    \newpage

    \question I think \underline{\hspace{2cm}} is the most serious crime. Killing someone is unacceptable.
    
    \begin{choices}
        \choice pickpocketing
        \choice vandalism
        \choice theft
        \CorrectChoice murder
    
    \end{choices}
    

    \question We have the same taste in music, so we always \underline{\hspace{2cm}} which bands are the best.
    
    \begin{choices}
        \choice worry about
        \choice wait for
        \choice argue with
        \CorrectChoice agree about
    
    \end{choices}
    
    
    \question This soup tastes horrible - I should have used a better \underline{\hspace{2cm}}.
    
    \begin{choices}
        \choice article
        \CorrectChoice recipe
        \choice advertisement
        \choice receipt
    
    \end{choices}
    

    	\section{Saberes y pensamiento científico}
\subsection{Matemáticas}

\question Si las raíces de la ecuación cuadrática $x^2 - 7x + k = 0$ son $r_1$ y $r_2$, y se sabe que $r_1^2 + r_2^2 = 29$, ¿cuál es el valor de la constante $k$?
    \begin{choices}
        \choice -10
        \choice 5
        \CorrectChoice 10 
        \choice 14
    \end{choices}

    \question En un triángulo rectángulo ABC, el ángulo recto está en C. Si el seno del ángulo A es igual a $3/5$ y la longitud de la hipotenusa AB es de 20 cm, ¿cuál es el perímetro del triángulo?
    \begin{choices}
        \choice 36 cm
        \choice 40 cm
        \CorrectChoice 48 cm 
        \choice 60 cm
    \end{choices}

    \question ¿De cuántas formas distintas se pueden ordenar las letras de la palabra "OLIMPIADA", garantizando que las tres "A" aparezcan siempre juntas?

    \begin{choices}
        \choice 1260
        \choice 5040
        \CorrectChoice 2520
        \choice 30240
    \end{choices}

    \newpage

    \question Un edificio proyecta una sombra de 18 metros en el mismo instante en que un poste vertical de 3 metros proyecta una sombra de 1.2 metros. Utilizando el teorema de Tales, la altura del edificio es de:
    \begin{choices}
        \choice 36 m
        \choice 42 m
        \CorrectChoice 45 m 
        \choice 54 m
    \end{choices}

    \question En una circunferencia de centro O, se traza un ángulo inscrito ABC que mide 42°. ¿Cuánto mide el ángulo central AOC que subtiende el mismo arco de circunferencia?
    \begin{choices}
        \choice 21°
        \choice 42°
        \choice 63°
        \CorrectChoice 84° 
    \end{choices}

      \subsection{Química}

      \question La organización original de la tabla periódica propuesta por Dmitri Mendeléyev predecía la existencia de elementos aún no descubiertos. ¿En qué parámetro físico basó principalmente su ordenamiento inicial?
    \begin{choices}
        \choice En el número atómico (cantidad de protones).
        \choice En la electronegatividad de Pauling.
        \CorrectChoice En las masas atómicas relativas y la repetición periódica de sus propiedades químicas.
        \choice En la configuración electrónica de los niveles de energía.
    \end{choices}

    \question Al analizar el enlace que forma el cloruro de sodio (NaCl), se determina que es un compuesto iónico debido a que:
    \begin{choices}
        \choice Los átomos de cloro y sodio comparten equitativamente sus electrones de valencia.
        \CorrectChoice Existe una transferencia completa de un electrón del sodio al cloro, formando iones de cargas opuestas que se atraen electrostáticamente.
        \choice Ambos átomos ceden protones al núcleo del elemento más electronegativo.
        \choice Se forma una red de núcleos positivos inmersos en un mar de electrones deslocalizados.
    \end{choices}

    \question ¿Cuáles son los coeficientes estequiométricos enteros más pequeños que balancean correctamente la ecuación de combustión del metano: $CH_4 + O_2 \rightarrow CO_2 + H_2O$?
    \begin{choices}
        \choice 1, 1, 1, 2
        \CorrectChoice 1, 2, 1, 2
        \choice 2, 3, 2, 4
        \choice 1, 3, 1, 2
    \end{choices}

    \question Los isótopos del carbono, como el Carbono-12 y el Carbono-14, son átomos del mismo elemento que difieren químicamente o físicamente en:
    \begin{choices}
        \choice La cantidad de protones en el núcleo atómico.
        \CorrectChoice La cantidad de neutrones en el núcleo atómico, lo que altera su masa.
        \choice La cantidad de electrones en el estado basal.
        \choice Su reactividad frente al oxígeno atmosférico.
    \end{choices}

    \question Durante una reacción química en un sistema cerrado, se comprueba la Ley de la Conservación de la Masa formulada por Lavoisier, la cual implica a nivel submicroscópico que:
    \begin{choices}
        \choice La energía potencial de los reactivos es igual a la de los productos.
        \choice El volumen total de los gases no cambia durante la reacción.
        \CorrectChoice Los átomos no se crean ni se destruyen, únicamente se reorganizan para formar nuevas moléculas.
        \choice Los enlaces químicos no se rompen, solo cambian de estado de agregación.
    \end{choices}
    
      \question La masa molar del ácido sulfúrico ($H_2SO_4$) es de aproximadamente 98 g/mol. Si un matraz contiene 49 gramos de ácido sulfúrico puro, ¿cuántas moléculas del ácido hay en el recipiente, considerando el número de Avogadro ($N_A = 6.022 \times 10^{23}$)?
    \begin{choices}
        \CorrectChoice $3.011 \times 10^{23}$ moléculas (0.5 moles).
        \choice $6.022 \times 10^{23}$ moléculas (1 mol).
        \choice $1.204 \times 10^{24}$ moléculas (2 moles).
        \choice $1.505 \times 10^{23}$ moléculas (0.25 moles).
    \end{choices}

    \question El pH es una escala logarítmica que mide la concentración de iones hidronio. Si el pH de una disolución cambia de 3 a 5, ¿qué le ocurre a la concentración de iones $H^+$ en la disolución?
    \begin{choices}
        \choice Se reduce a la mitad de su valor original.
        \choice Aumenta en un factor de 20.
        \CorrectChoice Disminuye en un factor de 100 veces (se hace más diluida).
        \choice Aumenta en un factor de 100 veces (se hace más concentrada).
    \end{choices}

    \question En una reacción de oxidación-reducción (redox), el agente oxidante es la especie química que:
    \begin{choices}
        \choice Cede electrones y aumenta su número de oxidación.
        \CorrectChoice Acepta electrones de otra sustancia, provocando la oxidación de esta última mientras él mismo se reduce.
        \choice Actúa como catalizador para acelerar la velocidad de la reacción.
        \choice Libera oxígeno diatómico al ambiente en estado gaseoso.
    \end{choices}

    \question Una reacción química cuyo diagrama de energía potencial muestra que la entalpía (energía potencial total) de los productos es menor que la entalpía de los reactivos se clasifica como:
    \begin{choices}
        \choice Reacción endotérmica, pues absorbe calor del entorno.
        \CorrectChoice Reacción exotérmica, pues libera energía térmica al entorno.
        \choice Reacción catalítica reversible en equilibrio estático.
        \choice Reacción isotérmica ideal sin intercambio energético.
    \end{choices}

    \question La química orgánica se fundamenta en la capacidad de los átomos de carbono de formar largas cadenas estables mediante enlaces covalentes. Esta propiedad excepcional se conoce como:
    \begin{choices}
        \choice Hibridación iónica.
        \CorrectChoice Concatenación.
        \choice Isomería óptica.
        \choice Electronegatividad extrema.
    \end{choices}

\section{De lo Humano a lo comunitario}
\subsection{Tecnología}
 	\question ¿Qué es una hoja de cálculo y para qué sirve?
    
    \begin{choices}
        \choice Es un procesador de textos diseñado para redactar cartas, memorándums y documentos formales de oficina.
        \CorrectChoice Es una herramienta digital que permite organizar, calcular y analizar información mediante tablas, filas y columnas, utilizada para operaciones matemáticas, registros y gráficas.
        \choice Es un programa de diseño vectorial que sirve primordialmente para la edición de imágenes, logotipos y maquetación web.
        \choice Es una base de datos relacional orientada únicamente al almacenamiento seguro de contraseñas y correos electrónicos.
    
    \end{choices}
    

    \question ¿Cuál es la diferencia principal entre Google Sheets y Microsoft Excel?
    
    \begin{choices}
        \choice Google Sheets es una herramienta de pago obligatorio, mientras que Microsoft Excel es de código abierto y completamente gratuito.
        \choice Microsoft Excel solo puede instalarse en sistemas operativos macOS, mientras que Google Sheets es exclusivo de Windows.
        \CorrectChoice Google Sheets funciona principalmente en línea y permite el trabajo colaborativo en tiempo real, mientras que Microsoft Excel es un programa más avanzado con funciones especializadas y uso óptimo sin conexión.
        \choice No existe ninguna diferencia técnica ni operativa entre ambas herramientas; comparten el mismo desarrollador y diseño.
    
    \end{choices}
    
    \question ¿Qué es una fórmula en una hoja de cálculo?
    
    \begin{choices}
        \CorrectChoice Es una operación que permite realizar cálculos automáticos y que siempre inicia con el signo igual (=), por ejemplo: \texttt{=A1+B1}.
        \choice Es una anotación informativa que se añade al pie de la página para documentar las fuentes bibliográficas de una tabla.
        \choice Es un comando estilístico utilizado exclusivamente para modificar el color de fondo, bordes y fuentes de un rango de celdas.
        \choice Es una secuencia de programación avanzada que sirve para borrar permanentemente los datos duplicados de un libro de trabajo.
    
    \end{choices}
    

    \question ¿Para qué sirve la función \texttt{SUMA} en Google Sheets y Excel?
    
    \begin{choices}
        \choice Para multiplicar de forma directa todos los valores numéricos contenidos en un rango seleccionado.
        \choice Para contabilizar cuántas celdas de un rango determinado contienen datos alfanuméricos o texto en lugar de números.
        \CorrectChoice Para agregar automáticamente varios números o celdas.
        \choice Para buscar de manera jerárquica un valor en la primera columna de una tabla y extraer la información correspondiente.
    
    \end{choices}
    
        \newpage


    \question ¿Qué ventajas tiene utilizar hojas de cálculo en actividades escolares o laborales?
    
    \begin{choices}
        \choice Permiten programar e implementar páginas web y aplicaciones móviles complejas sin requerir código.
        \CorrectChoice Ayudan a organizar información, realizar cálculos rápidos, crear gráficas, ahorrar tiempo y facilitar el análisis de datos de manera clara y ordenada.
        \choice Reemplazan la necesidad de usar sistemas operativos y optimizan automáticamente el ancho de banda de internet.
        \choice Sirven exclusivamente para el envío masivo de correos electrónicos corporativos y el manejo de redes sociales.
    
    \end{choices}
    

    \question ¿Qué son las funciones lógicas en Excel y Google Sheets?
    
    \begin{choices}
        \choice Son funciones matemáticas avanzadas empleadas para el cálculo automatizado de derivadas, integrales y matrices.
        \CorrectChoice Son herramientas que permiten tomar decisiones automáticas según una condición, ayudando a obtener resultados como ``verdadero'' o ``falso'' según los datos ingresados.
        \choice Son comandos mecánicos que ordenan de manera alfabética las filas de una tabla sin alterar las columnas adyacentes.
        \choice Son estilos de formato condicional que tiñen las celdas automáticamente dependiendo únicamente de si el número es par o impar.
    
    \end{choices}
    

    \question ¿Para qué sirve la función \texttt{SI} en Excel y Google Sheets?
    
    \begin{choices}
        \choice Sirve para sumar exclusivamente las celdas que cumplen con un criterio riguroso preestablecido por el usuario.
        \CorrectChoice Permite evaluar una condición y mostrar un resultado diferente según se cumpla o no.
        \choice Se utiliza para localizar una palabra clave específica dentro de un bloque extenso de texto en una sola celda.
        \choice Funciona para concatenar o enlazar el contenido textual de dos celdas independientes en una nueva columna.
    
    \end{choices}
    

    \question ¿Qué función cumplen las funciones \texttt{Y} y \texttt{O} en una hoja de cálculo?
    
    \begin{choices}
        \choice La función \texttt{Y} comprueba si al menos una condición es verdadera, mientras que la función \texttt{O} verifica que absolutamente todas se cumplan.
        \CorrectChoice La función \texttt{Y} verifica que todas las condiciones sean verdaderas, mientras que la función \texttt{O} comprueba si al menos una condición es verdadera.
        \choice Ambas funciones se emplean exclusivamente para realizar promedios ponderados y redondeos de cifras decimales complejas.
        \choice Sirven de puente de comunicación para enlazar la hoja de cálculo actual con sistemas gestores de bases de datos externos.
    
    \end{choices}
    
    \newpage

 \section{Ética, Naturaleza y Sociedades}
    \subsection{Formación Cívica y Ética}

    \question En el sistema democrático, el Plebiscito y el Referéndum son herramientas constitucionales cuyo objetivo es:
    \begin{choices}
        \choice Designar directamente a los ministros de la Suprema Corte de Justicia.
        \choice Auditar las cuentas bancarias de los partidos políticos en época electoral.
        \CorrectChoice Garantizar la democracia participativa, permitiendo a la ciudadanía consultar y votar sobre leyes o decisiones trascendentales de gobierno.
        \choice Sustituir al Poder Legislativo cuando entra en receso ordinario.
    \end{choices}

    \question El concepto de Desarrollo Sustentable, formalizado en el Informe Brundtland y esencial en la ética ambiental contemporánea, establece la necesidad de:
    \begin{choices}
        \choice Detener por completo el desarrollo tecnológico para regresar a métodos de producción artesanales.
        \choice Explotar rápidamente los recursos naturales en países en vías de desarrollo para igualar la economía mundial.
        \CorrectChoice Satisfacer las necesidades de las generaciones presentes sin comprometer la capacidad de las generaciones futuras para satisfacer sus propias necesidades.
        \choice Limitar el consumo de recursos de manera exclusiva a la población de los países industrializados.
    \end{choices}

    \question El Derecho Internacional Humanitario (DIH), resguardado por instrumentos como los Convenios de Ginebra, tiene como propósito ético y jurídico fundamental:
    \begin{choices}
        \choice Fomentar el libre mercado internacional eliminando los aranceles fronterizos.
        \CorrectChoice Limitar los efectos y la barbarie de los conflictos armados, protegiendo a civiles, prisioneros y personal médico.
        \choice Obligar a los países miembros a adoptar una moneda única a nivel mundial.
        \choice Establecer tribunales militares para juzgar crímenes comunes y robos internacionales.
    \end{choices}

    \question En el Estado de derecho, la cultura de la transparencia y la rendición de cuentas (garantizada en México por el INAI) funge como un antídoto social contra:
    \begin{choices}
        \choice La inflación y la devaluación cambiaria de la moneda nacional.
        \choice La migración documentada hacia el extranjero.
        \CorrectChoice La opacidad gubernamental, el tráfico de influencias y la corrupción institucional.
        \choice La participación ciudadana excesiva en las casillas electorales.
    \end{choices}

    \question El avance vertiginoso de la ciencia, como la manipulación genética, la clonación o el uso de inteligencia artificial en diagnóstico médico, demanda un análisis moral sistemático que proteja la dignidad humana. A esta rama de la ética se le denomina:
    \begin{choices}
        \choice Metafísica teleológica.
        \choice Ecología política.
        \choice Sociología del derecho.
        \CorrectChoice Bioética.
    \end{choices}

    \subsection{Historia}

\question A finales del siglo XIX y principios del XX, las potencias europeas intensificaron el Imperialismo. ¿Cuál fue la motivación económica subyacente que impulsó el reparto de África y Asia y detonó las tensiones previas a la Primera Guerra Mundial?
    \begin{choices}
        \choice La necesidad de importar mano de obra calificada hacia Europa.
        \CorrectChoice La búsqueda incesante de materias primas baratas y la creación de nuevos mercados cautivos para los productos industriales.
        \choice El objetivo humanitario de erradicar pandemias en el hemisferio sur.
        \choice La unificación de las iglesias cristianas en el continente asiático.
    \end{choices}

    \question En el marco de la Revolución Mexicana, el Plan de Ayala, promulgado por Emiliano Zapata en 1911, tuvo como demanda principal y central:
    \begin{choices}
        \choice La creación de un seguro social obligatorio para trabajadores de fábricas textiles.
        \choice El no a la reelección presidencial y el sufragio efectivo.
        \CorrectChoice La restitución inmediata de tierras, montes y aguas a los campesinos y pueblos originarios despojados por los hacendados.
        \choice La expulsión del clero extranjero del territorio nacional.
    \end{choices}

    \question Durante el periodo conocido como la Guerra Fría, ¿por qué los enfrentamientos entre Estados Unidos y la Unión Soviética se manifestaron en "guerras subsidiarias" (proxy wars) como las de Corea y Vietnam, en lugar de un conflicto bélico directo?
    \begin{choices}
        \choice Porque ambas potencias formaban parte del mismo consejo de seguridad nacional europeo.
        \CorrectChoice Por el temor a la Destrucción Mutua Asegurada (MAD) derivado del arsenal de armas nucleares de ambas superpotencias.
        \choice Porque las Naciones Unidas prohibieron terminantemente la guerra directa entre potencias del primer mundo.
        \choice Por la escasez de soldados disponibles tras las pérdidas de la Segunda Guerra Mundial.
    \end{choices}

    \question El periodo económico en México comprendido entre 1940 y 1970, caracterizado por altas tasas de crecimiento sostenido y baja inflación mediante la protección arancelaria, se conoce históricamente como:
    \begin{choices}
        \choice El Neoliberalismo Aperturista.
        \choice El Porfiriato Tardío.
        \CorrectChoice El Milagro Mexicano (Modelo de Sustitución de Importaciones y Desarrollo Estabilizador).
        \choice La Crisis de la Deuda Externa.
    \end{choices}

    \question La globalización contemporánea es un fenómeno multidimensional. A nivel cultural, uno de sus principales efectos y, a la vez, objeto de fuertes críticas sociales es:
    \begin{choices}
        \choice El aislamiento absoluto de las tribus urbanas.
        \CorrectChoice La homogeneización cultural y el riesgo de pérdida de tradiciones y lenguas locales frente al consumismo transnacional.
        \choice El cese inmediato de los flujos migratorios internacionales.
        \choice La fragmentación de los mercados bursátiles mundiales.
    \end{choices}
\end{questions}
\end{document}
